\documentclass[11pt,a4paper]{article}
% --- Packages ---
\usepackage[margin=1in]{geometry}
\usepackage{times}
\usepackage{amsmath}
\usepackage{graphicx}
\usepackage{enumitem}
\usepackage{titlesec}
\usepackage[colorlinks=true,linkcolor=black,citecolor=blue,urlcolor=blue]{hyperref}
\usepackage{xcolor}
\usepackage{fancyhdr}
\usepackage{booktabs}
% --- Section styling ---
\titleformat{\section}{\large\bfseries}{\thesection}{1em}{}
\titleformat{\subsection}{\normalsize\bfseries}{\thesubsection}{1em}{}
\titlespacing*{\section}{0pt}{14pt}{6pt}
\titlespacing*{\subsection}{0pt}{10pt}{4pt}
\pagestyle{fancy}
\fancyhf{}
\fancyhead[L]{\small Digital Image Processing - Course Project Proposal}
\fancyhead[R]{\small \thepage}
\renewcommand{\headrulewidth}{0.4pt}
\setlist[itemize]{topsep=2pt,itemsep=2pt,parsep=0pt}
\setlist[enumerate]{topsep=2pt,itemsep=2pt,parsep=0pt}
\begin{document}
% ================= TITLE =================
\begin{titlepage}
\centering
\vspace*{2cm}

{\LARGE \bfseries Seeing Through the Shake: Text-Aware PSF Estimation \\[4pt] for Blind Motion Deblurring of License Plates \\[4pt] and Low-Texture Text}

\vspace{1.2cm}
{\large Course Project Proposal}

\vspace{0.4cm}
{\large Digital Image Processing}

\vspace{2cm}
{\large
\begin{tabular}{c}
\textbf{Submitted by} \\[4pt]
Inesh Roy (2024112015) \\
Maithily Gauravkumar Bhala (2024117011) \\
\end{tabular}
}

\vspace{1cm}

\vfill

\end{titlepage}

% ================= 1. PROBLEM STATEMENT =================
\section{Problem Statement}

Motion blur is one of the most common forms of image degradation encountered outside
controlled laboratory conditions: a hand-held camera shakes during exposure, or a license
plate moves relative to a traffic camera during capture. Classical blind deblurring recovers
a sharp image without knowing the blur kernel in advance, by estimating the motion's
angle and length directly from the image's frequency-domain signature - typically via a
Radon transform on the Fourier spectrum (for angle) and a cepstral analysis (for length) -
and then inverting the blur with a non-blind deconvolution step such as Richardson-Lucy.
This approach is fully interpretable, requires no training data, and needs no learned model,
which makes it a natural fit for a Digital Image Processing project.

However, this pipeline was developed and is almost always evaluated on natural,
texture-rich photographs, where strong edges and varied structure give the Radon and
cepstral estimators a clean signal to lock onto. License plates, road signs, and printed text
are a fundamentally different class of input: small, low-contrast, and dominated by a
handful of straight character strokes rather than rich natural texture. This is precisely the
regime where legibility matters most and where blur is most damaging, yet it is also the
regime the classical pipeline is least tested on and, we hypothesize, weakest in - because
the sparse, directional structure of text can be too weak a signal for angle and length
estimators tuned on natural images, and because an incorrect PSF estimate does not merely
fail to help but can actively degrade the image further. Closing this specific gap - making
classical PSF estimation reliable on low-texture text - is the problem this project addresses.

% ================= 2. PROPOSED SOLUTION =================
\section{Proposed Solution}

We propose a text-aware blind deblurring pipeline built specifically to close the
natural-image-vs.-text performance gap described above, rather than a generic
reimplementation of the standard recipe.

\subsection{Core pipeline}
The backbone of the system follows the established classical approach: the log-magnitude
Fourier spectrum of the blurred image is analyzed with a Radon transform to recover the
blur angle, the 2-D cepstrum along that direction is used to recover the blur length, and the
resulting line-kernel PSF is inverted using Richardson-Lucy deconvolution (with Wiener
deconvolution implemented as a secondary comparison point).

\subsection{Text-contrast pre-processing}
Because text and plate regions offer the estimators a weaker directional signal than natural
photographs, we introduce a pre-processing stage applied before Radon/cepstrum analysis
that enhances character-stroke edges relative to background texture - operating on an
edge-enhanced, locally contrast-normalized representation of the image rather than raw
grayscale. The explicit goal of this stage is not general image enhancement but improving
the strength of the exact directional and periodic cues the angle and length estimators
depend on, specifically for this weak-signal input class.

\subsection{Confidence-gated restoration}
Since an unreliable PSF estimate can actively worsen the image rather than simply fail to
sharpen it, the pipeline includes a confidence check based on the strength of the detected
cepstral peak relative to background noise. When this confidence falls below a threshold,
the system withholds deconvolution and reports low confidence instead of applying an
unreliable PSF - making restoration a decision the pipeline makes deliberately, not a step
applied unconditionally.

% ================= 2.4 TEST DATA AND SCOPE =================
\subsection{Test data}
The pipeline is validated on three tiers of data, chosen specifically to probe the
natural-image-vs.-text gap: (i) synthetically blurred license-plate and printed-text crops
with known ground-truth angle and length, used to measure estimation error directly;
(ii) the Köhler camera-shake benchmark, used as an established real-blur reference point;
and (iii) real-world license-plate images from a public license-plate dataset, used for the
downstream OCR evaluation. We assume a single, spatially uniform blur kernel per image
throughout. The Köhler benchmark is used only as a global-kernel reference point for real
camera-shake blur, not to model spatial variation. For the natural-image contrast case in
Objective~2, we use synthetically blurred crops from a public natural-image set (e.g.\
Kodak or BSD), not the standard test images (Cameraman, Lena, House, Boat) used
elsewhere.

\subsection{Out of scope}
To keep the project distinct from a parallel classmate project tackling the same classical
pipeline, the following are explicitly out of scope here: (i) spatially-varying blur and
patch-wise PSF estimation; (ii) Wiener deconvolution as the primary restoration method
(we use Richardson-Lucy as primary and Wiener only as a secondary comparison); and
(iii) a head-to-head benchmark of Radon-based versus cepstrum-based estimators, which
we treat as one combined pipeline. That project instead addresses non-uniform,
patch-wise blur restoration on natural images with Wiener deconvolution as its primary
method, and compares the two estimators directly; we refer the reader there for that
setting.

% ================= 3. OBJECTIVES =================
\section{Objectives}

\begin{enumerate}
    \item Implement the core classical pipeline - Radon-based blur-angle estimation,
    cepstrum-based blur-length estimation, and Richardson-Lucy non-blind
    deconvolution - and validate it against known ground-truth kernels using
    synthetically blurred images and the Köhler camera-shake benchmark.
    \item Quantify, with a controlled sweep across blur angle, blur length, and noise level
    on uniform synthetic blur, the performance gap between natural, texture-rich images
    and low-texture, text-dominated images (license plates and printed text) for this
    classical pipeline.
    \item Design a text-contrast pre-processing stage targeted at strengthening the
    directional/periodic signal available to the Radon and cepstral estimators on
    low-texture inputs, and measure its effect against the same sweep used in
    Objective~2.
    \item Implement a confidence-gating mechanism that withholds deconvolution when
    the estimated PSF is judged unreliable, and evaluate how often it correctly avoids
    a harmful restoration versus how often it unnecessarily withholds a good one.
    \item Evaluate the complete text-aware pipeline on real-world license-plate images
    from a public dataset, using OCR accuracy before vs.\ after restoration as the
    primary evaluation criterion, and compare against the non-text-aware baseline
    from Objective~1.
\end{enumerate}

% ================= 4. DELIVERABLES =================
\section{Deliverables}

\begin{itemize}
    \item \textbf{Working implementation} of the full text-aware pipeline (PSF
    estimation, text-contrast pre-processing, confidence gating, non-blind
    deconvolution), implemented in Python using standard scientific-computing
    libraries.
    \item \textbf{Quantitative evaluation report} comparing the baseline classical
    pipeline against the text-aware pipeline, centered on OCR accuracy on real
    license-plate images before and after restoration, with angle/length estimation
    error and restored PSNR/SSIM on synthetic text/plate crops and the Köhler
    benchmark reported as supporting metrics.
    \item \textbf{Qualitative results gallery} of blurred vs.\ restored image pairs across a
    representative range of blur conditions, including documented cases where
    confidence gating correctly withholds restoration.
    \item \textbf{Project codebase} with a reproducible evaluation script.
    \item \textbf{Final presentation and report} summarizing the method, evaluation, and
    findings.
\end{itemize}

% ================= REFERENCES =================
\begin{thebibliography}{99}

\bibitem{fergus2006}
R.~Fergus, B.~Singh, A.~Hertzmann, S.~T.~Roweis, and W.~T.~Freeman,
``Removing Camera Shake from a Single Photograph,''
\emph{ACM Transactions on Graphics (Proc.\ SIGGRAPH)}, vol.~25, no.~3, pp.~787-794, 2006.

\bibitem{levin2009}
A.~Levin, Y.~Weiss, F.~Durand, and W.~T.~Freeman,
``Understanding and Evaluating Blind Deconvolution Algorithms,''
in \emph{Proc.\ IEEE Conference on Computer Vision and Pattern Recognition (CVPR)}, 2009.

\bibitem{radon-license-plate}
Y.~Wang et al.,
``A New Method and Implementation of Blind Restoration Algorithm for Moving Fuzzy
License Plate Image Based on Frequency-Domain Characteristics,''
\emph{IEEE Access}, 2021.

\bibitem{radon-blur-improved}
X.~Zhou, Y.~Zhang, and W.~Zhou,
``Improved Estimation of Motion Blur Parameters for Restoration from a Single Image,''
\emph{PLOS ONE}, 2020.

\bibitem{cepstrum-bitplane}
V.~Deshpande and M.~Patnaik,
``A Novel Modified Cepstral Based Technique for Blind Estimation of Motion Blur,''
\emph{Optik}, vol.~124, no.~24, pp.~6994-7000, 2013.

\bibitem{richardson1972}
W.~H.~Richardson,
``Bayesian-Based Iterative Method of Image Restoration,''
\emph{Journal of the Optical Society of America}, vol.~62, no.~1, pp.~55-59, 1972.

\bibitem{lucy1974}
L.~B.~Lucy,
``An Iterative Technique for the Rectification of Observed Distributions,''
\emph{The Astronomical Journal}, vol.~79, no.~6, pp.~745-754, 1974.

\bibitem{lpblur2024}
D.~Zhang et al.,
``A Dataset and Model for Realistic License Plate Deblurring,''
\emph{arXiv preprint arXiv:2404.13677}, 2024.

\bibitem{ccpd2018}
Z.~Xu et al.,
``Towards End-to-End License Plate Detection and Recognition: A Large Dataset and
Baseline,''
in \emph{Proc.\ European Conference on Computer Vision (ECCV)}, 2018.

\bibitem{kohler2012}
R.~Köhler, M.~Hirsch, B.~Mohler, B.~Schölkopf, and S.~Harmeling,
``Recording and Playback of Camera Shake: Benchmarking Blind Deconvolution with a
Real-World Database,''
in \emph{Proc.\ European Conference on Computer Vision (ECCV)}, 2012.

\end{thebibliography}

\end{document}